% ============================================================================
\section{The device concept: a ``?''-shaped ear-hook}
\label{sec:device}

Instead of a temple pod, the user proposed an ear device shaped like a question mark, worn on both ears. Its top curve carries dry EEG electrodes above and behind the ear (near T7/T8 and P7/P8), the back stem a mastoid reference, the ear-canal tip in-ear EEG, PPG and temperature, with an optional earlobe PPG, a 6-axis IMU, jaw-muscle (temporalis) EMG and Bluetooth to a phone.

{\footnotesize
\begin{verbatim}
        ___
      /     \   <- top curve: dry EEG electrodes above/behind the ear (~T7/T8, P7/P8)
     |       |
      \      |  <- back stem: mastoid reference electrode, temporalis EMG
       ?     |
       |    /   <- ear-canal tip: in-ear EEG + PPG + SpO2 + temperature
       .   /    <- optional earlobe PPG
   + 6-axis IMU (accelerometer, gyroscope) . Bluetooth to phone . worn on BOTH ears
\end{verbatim}
}

\begin{table}[H]
\centering\small
\caption{Why the ear.}
\begin{tabularx}{\linewidth}{@{}l L@{}}
\toprule
Reason & Evidence \\
\midrule
Over the temporal lobe, the commonest seizure origin in adults & Ear-EEG matched scalp EEG for temporal-lobe seizures (Zibrandtsen 2017) \\
Proven for seizure detection & Behind-the-ear 2-channel wearable caught 98\% of absence seizures (Swinnen 2021) \\
Good PPG site & Earlobe PPG up to 6.3$\times$ better signal than the wrist; in-ear SpO\textsubscript{2} valid 70--100\% \\
Discreet and stable & Looks like a hearing aid or earbud \\
Both ears $>$ one ear & 3--6$\times$ fewer false alarms than one-sided montages (CHB-MIT, Table~\ref{tab:chb3}); +6\% seizures and $-$25--30\% false alarms in the general model \\
\bottomrule
\end{tabularx}
\end{table}

When asked which configuration would be most accurate, the answer from Table~\ref{tab:generalft} was that the full cap with fine-tuning has the fewest false alarms (88/day) but that the two-ear hook catches about as many seizures, and that false alarms should be handled by additional signals (multimodal agreement, artefact sensors, tiered alerts, personal learning, risk forecasting), which led to the body-signal review and the SeizeIT2 study.

% ============================================================================
\section{The Epilepsy Live Lab: an offline interactive visualisation}
\label{sec:livelab}

The user proposed an HTML page with ``a field of glowing orbs of different colours'' fluctuating with the data, plus live graphs. The result, \code{viz/epilepsy_live.html}, is a single self-contained offline page (no libraries, no internet) fed by \code{viz/data.js} ($\approx$3.4\,MB), generated by \code{viz/export_data.py} from the saved ionic simulations, the inference results, the CHB-MIT data and the detector outputs.

\subsection{Layout and panels}

A one-screen, no-scroll grid (first 12$\times$6, later 12$\times$7) holds: the neuron field; live readouts; nine model charts (LFP, firing rate, \Ko{} with the Kalman-filter and naive device estimates, \Cli{} with the inferred value and \EGABA, \Nai{} and pump flux, adenosine, CBF/HbO/HbR, reversal potentials, and a 100-cell $\times$ 20\,s \Ko{} space map); an EEG panel; the detector probability; band powers with line length; and later a row of body sensors. Controls select the tissue (epileptic/healthy), playback speed, pause, and the EEG source. Deep links reproduce any view: \code{?scen=&t=&eeg=&speed=&paused=&type=&walk=&view=&h=&boot=}.

\subsection{The orb field (first version)}

The 100 neurons sat on a ring (excitatory blue-white, inhibitory pink), flashing on spikes; orange particles around each neuron represented its local \Ko; a gold core glow showed \Nai; cyan satellites orbiting each neuron showed \Cli{} loading; a violet haze showed adenosine; glowing blood vessels were coloured by the HbO/HbR balance and pulsed at the heart rate.

\subsection{Real vs model EEG and the virtual ear-EEG forward model}

The first EEG panel showed a 3-minute, 5-channel excerpt of CHB-MIT \code{chb01_03.edf} (2926--3106\,s file time; clinician seizure at 2996--3036\,s), with the trained both-temples detector. The user asked whether model and real EEG were the same source (they were not: separate clocks), and proposed replacing the real recording with ``recordings from the model''. A forward model was therefore written that computes what ear electrodes around the simulated patch would record:
\begin{equation}
\mathrm{ch}_e(t) = g\sum_c w_{ec}\,s_c\,(\mathrm{spikes}_c * \kappa)(t) \;-\; \alpha\cdot51\,\mathrm{mV}\cdot\log_{10}\frac{[\mathrm K]_{o,\mathrm{local}}}{[\mathrm K]_{o,\mathrm{rest}}} \;+\; \eta_{1/f}(t),
\end{equation}
with Gaussian ring-distance weights $w_{ec}$ ($\sigma=20$ cells), signs $s_c=+1$ (excitatory) or $-1.5$ (inhibitory), a synaptic kernel $\kappa(t)=e^{-t/8\,\mathrm{ms}}-e^{-t/2\,\mathrm{ms}}$, a gain $g$ set so the healthy baseline carries $\approx$10\,\uV{} RMS of neural signal (shared between tissues), a glial \Kp{} DC shift reaching the scalp with $\alpha=1$\%, and pink background noise (6\,\uV, half shared across channels; 3\,\uV{} on the DC channel). Four virtual electrodes (V1--V2 right ear, V3--V4 left ear) plus a DC channel were exported at 128\,Hz with the same 14 features as the real detector.

A line-length detector \textbf{failed} on this virtual EEG because the model's ictal event is tonic firing, which band-passed EEG barely separates from interictal bursts. The detector was replaced by a \Kp-DC detector: the causal 1\,s mean of the DC channel below $-150$\,\uV{} (local \Ko{} $\gtrsim$8\,mM) raises an alarm, with a sigmoid probability $1/(1+e^{(\mathrm{DC}+150)/15})$ and a 60\,s refractory period. The threshold was chosen on the model itself and is not independently validated. In epileptic tissue it alarms at $t=11$\,s (1\,s after trigger onset); in healthy tissue it stays silent.

\subsection{Body sensors and EEG-only vs multimodal detection}

To show how body signals affect detection, a body module generates the SeizeIT2 modalities from the simulated seizure. Because the 100-neuron patch models only the onset zone, the spread to autonomic and motor areas uses literature timings (clinical seizure 40\,s from the trigger; ictal HR +53\% in focal seizures, $\approx$+100\% in tonic-clonic). Four seizure types are selectable: focal non-motor, focal motor (hyperkinetic), focal-to-bilateral tonic-clonic, and ``no HR response'' ($\approx$18\% of patients). A 12\,s walk episode can be switched on.

\paragraph{Heart.} A sympathetic drive $s$ follows the clinical seizure with rise/recovery time constants 6/50\,s; walking adds a drive $w$ (8/20\,s). Heart rate on a 10\,Hz grid is $HR = HR_0(1+h_{type}s) + 18w + d$ with $HR_0=68$\,bpm and a slow random drift $d$. Beats are generated with respiratory sinus arrhythmia at 0.25\,Hz whose amplitude (3\,bpm) shrinks with vagal withdrawal $(1-0.8s)$, plus 10\,ms jitter. The ECG is synthesised as Gaussian P, Q, R, S, T waves per beat (ECGSYN-style), with the T wave placed by Bazett $\sqrt{RR}$, baseline wander, noise and muscle contamination.

\paragraph{Muscle and motion.} Shoulder EMG is white noise with an envelope: 5\,\uV{} baseline; arm swing when walking; irregular 45\,\uV{} activity in hyperkinetic seizures; in tonic-clonic seizures a continuous 130\,\uV{} tonic phase (10\,s), then clonic bursts of 120\,ms slowing from 4 to 1\,Hz, then postictal atonia. Accelerometer magnitude shows 1.8\,Hz step impulses when walking, irregular large movements in hyperkinetic seizures, a 9\,Hz tremor in the tonic phase and decaying jerk impulses in the clonic phase.

\paragraph{Walking artefact.} Each step rocks the ear electrodes (biphasic transients) and a slow electrode-pressure drift of $-240$\,\uV{} appears on the DC channel, which deliberately fools the EEG-only \Kp-DC detector.

\paragraph{Fusion rule (illustrative, hand-set).} On a 1\,s grid, $HR_{10}$ is the mean instantaneous HR over the last 10\,s, a causal 1-minute baseline $B$ tracks it, activity $a$ is the 10\,s standard deviation of acceleration, and the heart rate the movement explains is $E = B + \min(20, 150a)$. EMG ratio $r$ is the 4\,s RMS over its causal 2-minute log-baseline. Body evidence is
\begin{equation}
b = \max\!\Big(\sigma\big((HR_{10}-E-12)/2.5\big),\ \sigma\big((r-4)/0.6\big)\Big),
\end{equation}
and the multimodal score multiplies the maximum EEG probability over the last 30\,s by the maximum body evidence over the last 10\,s. An early version looked back 40\,s for body evidence; when a walk followed a seizure, lingering postictal tachycardia combined with the walking artefact produced a false multimodal alarm, so the window was shortened to 10\,s (``current'' body evidence).

\begin{table}[H]
\centering\small
\caption{Live Lab results with the walk episode on.}
\begin{tabular}{@{}lll@{}}
\toprule
Body setting & EEG-only & Multimodal (EEG confirmed by body) \\
\midrule
Focal non-motor & caught +1 s, 1 false alarm (walk) & caught +9 s, 0 false \\
Focal motor & caught +1 s, 1 false & caught +5 s, 0 false \\
Tonic-clonic & caught +1 s, 1 false & caught +6 s, 0 false \\
No HR response & caught +1 s, 1 false & \textbf{missed}, 0 false \\
Healthy tissue & 1 false & 0 false \\
\bottomrule
\end{tabular}
\end{table}

This anticipated, qualitatively, the trade-off later measured on real data: body confirmation removes movement false alarms but costs a few seconds and misses seizures without a body sign. The live score first printed ``MISSED'' before the confirmation window had elapsed; it now shows ``watching\dots'' during a seizure.

\subsection{The 3D neuron plane}

The user asked for something ``more insane'': a plane of neurons fluctuating on the $z$ axis with fixed $x,y$. The 100 neurons were laid out on a 10$\times$10 grid in ring order snaking row by row, so ring neighbours stay plane neighbours. Each neuron's height is its \emph{own} data: local \Ko{} (slow swell), a spike trace (fast jolt), or both (selectable). A surface interpolated between neurons (visual only) is coloured by \Kp{} with a magma palette and depth-sorted; spikes emit ripples; orange \Kp{} ``vapour'' rises in proportion to local \Kp{} excess; \Clm{} satellites orbit; blood vessels run beneath; the view auto-rotates and supports drag, zoom and reset, all in plain canvas 3D projection. Initial renders were too large and blew out to white at the seizure peak; heights, perspective, mote rate and glow were reduced. The original ring remains one click away. Only 100 neurons exist in the simulation; a denser plane would require re-running \code{ionic.py} with more cells.

\subsection{The HUD interface}

Finally the interface was restyled as a ``Stark-lab'' heads-up display: a boot sequence (spinning arc reactor and typed log lines, skippable with \code{?boot=0}); holographic panels with corner brackets, a grid backdrop, scanlines and a sweeping beam; arc-reactor ring gauges for the eight readouts (conic gradients whose cores brighten with the value); glowing chart traces; a targeting reticle over the neuron field whose colour shifts from cyan to red with \Kp{} load and which flashes ``ICTAL THRESHOLD''; a flashing alert badge in ictal phases; and a ``JARVIS'' narration box. Everything stays offline (system fonts only).

\subsection{Bugs fixed in the Live Lab}

JSON serialisation of int64 values failed in the exporter; a CSS class \code{dash} for dashed legend keys clashed with the \code{.dash} grid and was renamed \code{dsh}; a negative \code{requestAnimationFrame} time step produced NaN and was clamped at zero; the multimodal 40\,s look-back and premature ``MISSED'' label were fixed as described; and the headless-browser screenshot job used for checking hung on the endless animation until a hard timeout was used.
